This discussion section is interpretive.
It is not part of the formal proof itself.

\paragraph{Purpose of the current revision.}
The manuscript now presents one short proof spine in the main text:
(1) define one canonical balance predicate $\Phi_H(s)=0$,
(2) state correspondence,
(3) state one canonical exclusion target,
(4) derive RH conditionally.
This is a presentation and scoping decision for audit clarity.
It does not reduce the technical burden of the open steps.

\paragraph{What the framework attempts to show.}
The argument attempts to show that if the canonical balance predicate is put in exact correspondence with nontrivial zeros, and if the canonical exclusion statement
\[
\mathcal{B}\subseteq\{s\in\mathcal{A}_{\mathrm{strip}}:\RePart(s)=\tfrac12\}
\]
is proved, then RH follows.
The framework suggests a spectral-coherence route to those two targets, but those targets remain open in the current manuscript version.

\paragraph{Why mathematical branching was reduced in the main flow.}
Earlier drafts carried multiple near-equivalent predicates and exclusion formulations in parallel.
Those variants are useful for technical experimentation, but they obscure claim traceability in the main proof narrative.
In this revision, the main text keeps one canonical predicate and one canonical exclusion statement; alternatives are treated as supporting material only.

\paragraph{Current mathematical status.}
The manuscript remains a proposed proof framework under review.
The RH deduction in the main text is conditional on open theorem targets.
No claim of complete proof is made.

\paragraph{Interpretive boundary.}
Heuristic motivation, computational checks, and architectural planning can guide research, but they are not substitutes for theorem-level proofs.
Formal validity depends only on statements proved under explicit hypotheses.
